\documentclass[10pt,a4paper]{article}
\usepackage{amsmath,amssymb}

\title{Lab Report: Damped Oscillations of a Pendulum}
\author{Chloe Baptiste \and Jonas Kral \and Aiko Tanaka}
\date{November 22, 2017}

\begin{document}
\maketitle

\section{Purpose}
The aim of this experiment was to measure the damping constant of a simple pendulum and to compare it with the prediction of a linear model.

\section{Method}
We used the following equipment:
\begin{description}
  \item[Pendulum] a brass bob of mass $0.25$~kg on a thin string.
  \item[Timer] a photogate, accurate to one millisecond.
  \item[Protractor] used to set the initial angle of $5^\circ$.
\end{description}
The procedure was as follows.
\begin{enumerate}
  \item Release the bob from the initial angle.
  \item Record the amplitude after each full swing.
  \item Repeat the measurement ten times and average.
\end{enumerate}

\section{Theory}
For small angles the angle $\theta(t)$ obeys
\begin{equation}
  \ddot\theta + 2\beta\dot\theta + \omega_0^2\theta = 0,
\end{equation}
whose underdamped solution is $\theta(t)=\theta_0 e^{-\beta t}\cos(\omega t+\varphi)$ with $\omega=\sqrt{\omega_0^2-\beta^2}$. The period is $T=2\pi/\omega$, and the amplitude after $n$ swings is $\theta_n=\theta_0e^{-\beta nT}$.

\section{Results}
\begin{table}[htbp]
\centering
\begin{tabular}{lrr}
\hline
Swing $n$ & Amplitude (deg) & $\ln(\theta_n/\theta_0)$\\
\hline
0 & 5.00 & 0.000\\
5 & 4.12 & $-0.194$\\
10 & 3.39 & $-0.389$\\
20 & 2.31 & $-0.772$\\
\hline
\end{tabular}
\caption{Measured amplitudes.}
\label{tab:amp}
\end{table}

A straight-line fit of the last column against $n$ in Table~\ref{tab:amp} gives a slope of $-0.0386$ per swing. With $T=1.84$~s this yields $\beta=0.021\ \mathrm{s}^{-1}$, with an uncertainty of about $5\%$.

\section{Discussion}
The agreement with the linear model is good; the small curvature seen at large $n$ may come from air resistance that grows with speed.\footnote{A quadratic drag term would bend the semilog plot.} The main source of error was reading the protractor.

\end{document}
